\documentclass[11pt,a4paper]{article}
\usepackage[margin=1in]{geometry}
\usepackage[round,authoryear]{natbib}
\usepackage{setspace}
\usepackage[colorlinks=true,linkcolor=black,citecolor=blue,urlcolor=blue]{hyperref}
\usepackage{graphicx}
\usepackage{amsmath}
\usepackage{amssymb}
\usepackage{titling}

% Springer-style keywords environment
\providecommand{\keywords}[1]{\vspace{0.5em}\par\noindent\textbf{Keywords} \quad #1}

\title{Recursive Generative-Adversarial Delegation:\\[0.3em]
\large An Interdisciplinary Synthesis for Human-Governed Multi-Agent AI Systems}

\author{%
  Ageless Kao\thanks{Corresponding author: \texttt{ageless\_kao@trendmicro.com}. The views expressed in this paper are those of the author and do not necessarily reflect those of TrendAI or Trend Micro Inc.}\\
  \small Group Product Manager, TrendAI, Trend Micro Inc.\\
  \small Ottawa, Canada
}

\date{\small Submitted: \today}

\begin{document}
\maketitle

\begin{abstract}
Multi-agent AI systems increasingly take the form of recursive delegation hierarchies under human oversight---a topology already developed in iterated amplification and factored cognition \citep{christiano2018ida, christiano2016hch}, hierarchical multi-agent system taxonomies \citep{moore2025}, and production manager-worker frameworks such as MetaGPT and CrewAI \citep{hong2024metagpt}. This paper proposes the Recursive Generative-Adversarial Delegation (RGAD) framework not as a structural discovery but as an interdisciplinary synthesis lens over these architectures, integrating Generative Adversarial Networks, Hegelian dialectics, AI debate, actor-critic LLM patterns, and human--machine communication theory.

RGAD analyzes the role-space along a single dyadic axis---Generator (produces, commits) and Validator (challenges, discriminates)---with internal nodes occupying both roles simultaneously and the human at the apex as terminal, non-delegable Validator. We treat this dyadic reading as a working proposal rather than a proven minimum.

The paper's central structural claim is what we call the \emph{Generator-Validator coupling thesis}: at any single agent, validation competence is functionally inseparable from sustained generative engagement, because the discriminative judgments a Validator must make depend on the active generative activity that knows what a good candidate looks like. We propose this thesis as a structural-causal generalization of older empirical and philosophical traditions---Goodall's Theory of Expert Leadership, Polanyi's tacit-knowledge framework, and the software-engineering code-review literature---extended to AI delegation in the new cost regime, where generative AI has collapsed the cost of generation while leaving validation roughly unchanged. Asymmetric delegation (outsourcing generation while nominally retaining validation) therefore atrophies discriminative competence specifically. Recent randomized evidence supports this: participants who sustained generative engagement preserved discriminative skill, while those who delegated generation showed pronounced deficits in debugging \citep{shen2026}. The coupling thesis relocates the critical infrastructure of human-governed AI from human authority to the \emph{engaged human mind}: the cognitive activity of sustained generative participation.

We further read persona-based agent architectures across creative collaboration, decision support, human--robot interaction, and explicit system-level frameworks such as ASAF as social-affordance scaffolding over the Generator/Validator basis, and situate these mechanisms within a civilizational trajectory in which humans increasingly occupy Validator roles.

\keywords{Multi-agent AI systems \and Recursive delegation \and Generator-Validator coupling \and Engaged human mind \and Cognitive scaffolding \and Human oversight \and AI deskilling \and Critical infrastructure}
\end{abstract}

\section{Introduction}
As large language models (LLMs) are embedded in increasingly complex multi-agent systems, attention has focused on orchestration mechanisms, tool integration, and token efficiency. Several literatures already articulate the basic architecture under which these systems operate: iterated amplification (IDA) is a training strategy that approximates the Humans Consulting HCH (HCH) construction, in which a human at the apex decomposes a task into sub-tasks delegated to sub-agents who may themselves recursively delegate \citep{christiano2018ida, christiano2016hch}; recent taxonomies of hierarchical multi-agent systems formalize this design space along multiple coordination axes and survey generator/evaluator-style reflection loops among the coordination patterns \citep{moore2025}; and production frameworks---MetaGPT, CrewAI, Google's Agent Development Kit, LangGraph---implement coordinator-worker hierarchies in which a manager agent decomposes goals, dispatches to specialist workers, and validates returned outputs \citep{hong2024metagpt}. At the single-task level, multi-agent LLM literature has converged on Actor-Critic and Maker-Checker patterns that separate generation from criticism \citep{yu2024acccollab}.

What is less developed is a unified epistemic reading of this widely-deployed topology that connects it to its dialectical, generative-adversarial, debate-theoretic, and human--machine-communication antecedents. This paper offers such a reading.

Two intellectual traditions, in particular, illuminate the structure that already exists. The Generative Adversarial Network (GAN) architecture couples a Generator with a Discriminator in a zero-sum game, converging toward an equilibrium where neither can improve without the other \citep{goodfellow2014}. Hegelian dialectical reasoning frames knowledge development as the tension and resolution between Thesis and Antithesis in a Synthesis. Both share a common structure: a generative pole that commits to candidate outputs and an adversarial pole that challenges, discriminates, and rejects. AI debate \citep{irving2018} and scalable oversight more broadly extend this dyad into safety-relevant configurations.

The Recursive Generative-Adversarial Delegation (RGAD) framework proposed here is an interdisciplinary synthesis lens over these strands. It does not claim to introduce the recursive-tree topology, the human-at-apex framing, or the generator/validator dyad as such; each is well-established in the prior literature surveyed above. What RGAD proposes is:

\begin{enumerate}
\item[(i)] A unified terminological treatment that reads the GAN Discriminator, the Hegelian Antithesis-producer, and the AI-debate judge as instances of a common epistemic primitive (``Validator'').
\item[(ii)] An explicit framing of the dual-occupancy of internal nodes (Generator upward, Validator-Orchestrator downward) as the same primitive operating in two directions.
\item[(iii)] A reading of recent persona-based agent architectures---persona-split LLMs in creative collaboration \citep{epperson2025}, conflicting-persona ensembles in decision support \citep{kumar2025}, personality-differentiated agents in human--robot interaction \citep{alami2026}, and explicit system-level frameworks such as ASAF \citep{lee2026}---as a social-affordance scaffolding layer over the Generator/Validator basis.
\item[(iv)] The \emph{Generator-Validator coupling thesis}: at any single agent, validation competence is functionally inseparable from sustained generative engagement, because the discriminative judgments a Validator must make depend on the active generative activity that knows what a good candidate looks like. We propose this thesis as a structural-causal generalization of two prior traditions---Goodall's empirical Theory of Expert Leadership \citep{goodall2011} and Polanyi's tacit-knowledge tradition \citep{polanyi1966}, with software-engineering-specific empirical support from the code-review literature \citep{bacchelli2013}---applied to AI delegation in the new cost regime created by generative AI. It has a sharp consequence for asymmetric delegation: outsourcing the Generator role to AI while retaining the Validator role nominally severs the coupling at the apex of the recursive tree, where validity terminates. The thesis predicts a specific empirical signature in the form of disproportionate atrophy of discriminative skills under AI assistance, consistent with recent randomized evidence on debugging-specific deficits \citep{shen2026}. Together with the qualitative finding from that evidence---that engagement pattern, not AI use itself, is decisive---the coupling thesis relocates the critical infrastructure of human-governed AI from human authority to the engaged human mind.
\end{enumerate}

The contribution is therefore best characterized as synthesis and reframing, plus the structural-causal generalization in (iv) that explains why the apex node is a single point of failure and predicts the empirical signature of its compromise. The framework's value, if any, lies in the integration of literatures that typically do not speak to one another, in the empirical and normative implications that follow from reading them together, and in the falsifiable prediction that the coupling thesis makes about which competences atrophy under AI delegation.

\section{Background and Related Work}

\subsection{Generator/Discriminator Dyad in Machine Learning}
The GAN architecture introduced by Goodfellow et al. established a powerful generative modeling paradigm in which a Generator produces samples and a Discriminator attempts to distinguish generated samples from real data \citep{goodfellow2014}. Optimal training corresponds to a Nash equilibrium in which neither player can improve without the other.

Multi-agent extensions such as MAD-GAN deploy multiple Generators against a single Discriminator to mitigate mode collapse and improve diversity \citep{ghosh2018}. Other work extends adversarial structures to additional domains, including tabular data synthesis and collaborative image editing. In each case, the generative/adversarial split remains the core structural primitive.

\subsection{Dialectical Reasoning and AI Debate}
Dialectical reasoning, systematized by Hegel, formalizes knowledge development as the tension and resolution between Thesis and Antithesis in a Synthesis. Recent work implements dialectical structures in LLMs for reasoning evaluation and self-improvement, including approaches that use multiple agents to produce thesis and antithesis positions and approaches that implement internal debate within a single LLM to refine outputs.

In AI safety, Irving et al.'s debate framework posits two AI debaters presenting competing arguments to a human judge, who acts as the final arbiter \citep{irving2018}. This structure---two generators and a human validator---anticipates the RGAD configuration at depth one.

\subsection{Recursive and Hierarchical Multi-Agent Architectures}
Several recent architectures allow agents to recursively spawn sub-agents and decompose tasks into hierarchies. THREAD and ReDel introduce frameworks for dynamic recursive delegation and flexible decomposition of complex problems \citep{zhu2024thread, zhu2024redel}.

These frameworks show that recursive delegation is practically feasible, but they leave the epistemic roles of agents implicit. RGAD makes the Generator/Validator structure explicit at each node, arguing that this is the necessary configuration for self-correcting collaboration.

\subsection{Human-in-the-Loop and Scalable Oversight}
Human-in-the-loop (HitL) paradigms distinguish between humans intervening during AI execution, supervising after completion, or being removed from the loop entirely. AI agents still fail in many multi-step tasks when unsupervised, underscoring the need for structured human oversight.

Scalable oversight research explores how weak human judges can reliably evaluate stronger AI systems. AI debate, recursive reward modeling, and self-critique methods all assume an ultimate human judge whose role is to discriminate between better and worse AI outputs \citep{irving2018}. Rahwan's Society-in-the-Loop framework generalizes this idea to social-level oversight structures \citep{rahwan2017}. In RGAD terms, these judges are instances of the Validator primitive.

\subsection{Social Affordances and Human--Machine Communication}
The RGAD and ASAF frameworks build directly on a longstanding literature in human--machine communication (HMC) and affordance theory. The Computers Are Social Actors (CASA) paradigm demonstrated that people spontaneously apply social scripts to computers and media technologies, responding to them as if they were social actors \citep{reeves1996}. Subsequent work has refined CASA, suggesting that while the effect may weaken for familiar, mundane devices, it remains strong for emergent, agentic, or anthropomorphic technologies. This confirms that social identity cues are a robust design lever for structuring human expectations of software systems.

Two adjacent threads from this tradition shape the analysis below. Suchman's account of human--machine reconfigurations argues that human--AI collaboration is best understood not as a fixed division of labor between cognitive agents but as a dynamic configuring and reconfiguring of human and machine roles in situated practice \citep{suchman2007}; this practice-oriented framing is consonant with the engaged-mind reading developed in Section 5.1, in which the human Validator's competence is constituted by ongoing generative engagement rather than by static role assignment. Affordance theory in communication studies provides a complementary structural lens: Hutchby argued that technologies possess communicative affordances---relational properties that both enable and constrain possible actions, grounded in material and design characteristics rather than purely in user perception \citep{hutchby2001}; Nagy and Neff extended this with the concept of imagined affordances, emphasizing that what users believe a system can do---based on design signals, prior experience, and cultural narratives---can be as behaviorally consequential as what the system is actually capable of \citep{nagy2015}.

Persona-based agent architectures, including ASAF \citep{lee2026} and the related frameworks surveyed below, can be understood as domain-specific operationalizations of imagined affordances for multi-agent LLM systems. Agent personas and role labels are designed identity signals that shape users' imagined affordances of each agent---what it is appropriate to ask, how much to trust, and how adversarially to engage. The RGAD analysis in this paper reinterprets these persona effects as a structured cognitive scaffolding layer over the Generator/Validator basis: a way of using social scripts and imagined affordances to make the underlying epistemic roles legible and navigable at scale.

Recent work on persona-based LLM systems underscores that this scaffolding pattern is emerging across multiple domains. Persona-split LLMs that separate divergent and convergent roles have been shown to scaffold human creative problem-solving by structuring when users should expand versus narrow the space of ideas \citep{epperson2025}. Multi-persona decision-support frameworks use deliberately conflicting agent personas to surface value tensions and extend human reflection, treating cognitive friction as a design feature rather than a defect \citep{kumar2025}. In human--robot interaction, multi-agent systems that give each robot a distinct personality and long-term memory improve users' ability to differentiate agents and personalize interactions, while central coordination reduces overlap and improves overall coherence \citep{alami2026}. Persona-based evaluator agents and personalized long-horizon assistants further demonstrate that personas can function as control structures guiding evaluation and action, not merely as stylistic choices \citep{li2025, zhang2025}. We read these as a system-level pattern: persona architectures use social identity as a cognitive scaffold for navigating complex agent ecosystems.

\subsection{Relationship to Prior Frameworks}
The architectural moves RGAD makes are not novel in isolation; what is offered is their integration. Four prior literatures already articulate components of the framework, and this subsection clarifies what RGAD adds and what it does not claim to add.

\textit{Iterated amplification and factored cognition.} Christiano's Humans Consulting HCH (HCH) construction posits a recursive tree of agents in which a human at the root decomposes a task into sub-tasks delegated to sub-agents, who may themselves recursively delegate; iterated amplification (IDA) is the training strategy designed to approximate this idealized process \citep{christiano2018ida, christiano2016hch}. The recursive-tree topology with a human at the apex---one of the architectural primitives we use---originates in this lineage. RGAD does not claim to introduce this topology. Our contribution is to make explicit the epistemic role each node plays: each internal node is not merely a delegator but simultaneously a Generator to its parent and a Validator to its children, with adversarial structure preserved at every depth.

\textit{Hierarchical multi-agent system taxonomies.} Moore's HMAS taxonomy organizes hierarchical multi-agent designs along five axes (control hierarchy, information flow, role/task delegation, temporal layering, communication structure) and explicitly discusses Generator-Evaluator reflection loops as a coordination pattern \citep{moore2025}. RGAD can be read as a particular configuration within this taxonomic space: static role hierarchy along the role-delegation axis, top-down control flow, and a Generator-Evaluator pattern instantiated at every depth. We do not propose RGAD as an alternative to this taxonomy; we propose it as an epistemically-explicit reading of one common pattern within it.

\textit{Production manager-worker frameworks.} Industrial multi-agent platforms---MetaGPT \citep{hong2024metagpt}, CrewAI's hierarchical process, Google's Agent Development Kit, and LangGraph's supervisor pattern---implement coordinator-worker hierarchies in which a manager agent decomposes goals, dispatches to specialist workers, and validates returned outputs. The internal node's dual-occupancy framing (Generator upward, Validator-Orchestrator downward) closely mirrors what these frameworks already do in practice. RGAD's value over them is not a new architecture but a unifying conceptual frame that connects this engineering pattern to alignment-theoretic and dialectical traditions.

\textit{Actor-Critic and Maker-Checker LLM patterns.} Multi-agent LLM literature has converged on numerous generator/validator dyads at the single-task level: Actor-Critic frameworks for collaborative reasoning \citep{yu2024acccollab}, Maker-Checker patterns that separate generation from criticism, and Generator-Critic-Refiner loops with iterative repair. RGAD generalizes these task-level patterns to a system-level topology under explicit human governance, and connects them to the Hegelian and GAN traditions through a unified terminological treatment.

\textit{Distributed cognition.} The dual-occupancy of internal nodes (Generator upward, Validator-Orchestrator downward) echoes Hutchins's distributed-cognition framework, in which cognitive activity is distributed across humans, artifacts, and social structures rather than localized in any single mind \citep{hutchins1995}. RGAD adds explicit adversarial structure at every depth, and---through the Generator-Validator coupling thesis developed in Section 5.1---identifies the engaged human mind at the apex coupling as a structural single point of failure that distributed-cognition accounts do not in themselves predict.

\textit{Expert leadership and tacit knowledge.} A long-standing tradition in management science and philosophy of knowledge documents the same coupling at the human--human level. Goodall's empirical Theory of Expert Leadership shows, across nearly two decades of cross-domain studies (hospitals led by physicians, universities led by scholars, sports teams led by former players), that organizations led by people who have done the core work themselves outperform those led by professional managers without that grounding \citep{goodall2011}. The hospital-CEO study finding a strong positive association between top-100 quality ranking and physician leadership is the most-cited result; the broader thesis is that expert leaders bring deep intuitive knowledge that shapes judgment in ways professional management training does not replicate. Polanyi's tacit-knowledge tradition provides the philosophical foundation: genuine competence rests on embodied knowledge that cannot be fully articulated and therefore cannot be transferred to a non-engaged validator \citep{polanyi1966}. In software engineering specifically, empirical research on code review documents that effective discriminative judgment requires sustained engagement with the codebase \citep{bacchelli2013}; the long-running practitioner debate over whether engineering managers should continue writing code is precisely this coupling question under a different name. The Generator-Validator coupling thesis developed in Section 5.1 is, at root, a structural-causal generalization of these traditions, applied to AI delegation in the new cost regime created by generative AI.

In summary, the structural elements of RGAD---recursive delegation, human at apex, generator/validator dyad, dual-occupancy at internal nodes---are each present in the prior literature. The contribution we offer is the interdisciplinary synthesis that reads them as instances of a common epistemic structure, the consolidated terminology (Validator as primitive across GAN, dialectics, and debate), the reading of persona-based agent design as cognitive scaffolding over this primitive, and the connection of architectural depth to skill-formation evidence on human Validator capacity.

\section{The Recursive Generative-Adversarial Delegation Framework}

\subsection{A Candidate Minimal Dyadic Basis}
\textbf{Claim 1 (proposed).} A candidate minimal basis for self-correcting collaborative reasoning is a single dyadic axis: Generator (produces, synthesizes, commits) and Validator (challenges, discriminates, rejects). We treat this as a working proposal rather than a proven minimum, motivated by structural convergence across the traditions reviewed below.

Throughout this paper, \emph{Validator} is used as the general term for the challenging, discriminating pole of this dyad. It is functionally equivalent to the Discriminator in GAN architecture, the Antithesis producer in Hegelian dialectics, and the judge in AI debate frameworks \citep{goodfellow2014, irving2018}. The term ``Discriminator'' appears only when directly referencing GAN literature; ``Validator'' is used in all other contexts, including when describing human roles, internal agent roles, and system-level positions in the RGAD tree.

This terminological unification reflects the core structural reading we propose: the GAN Discriminator, the Hegelian Antithesis producer, and the human oversight judge can be read as instances of a common epistemic primitive---an agent structurally committed to challenging and rejecting outputs from the generative pole. The mathematical formalisms differ; we suggest that the epistemic function is the same.

The proposal is motivated by the convergent structure of GAN theory, dialectical reasoning, and the debate-based scalable oversight literature \citep{goodfellow2014, irving2018}: across these traditions, systems that can both generate candidate outputs and validate between better and worse outputs exhibit the structural ingredients for self-improvement. We do not claim a formal proof of minimality; convergence across three frameworks is suggestive rather than conclusive. We therefore treat the dyad as a useful organizing primitive whose adequacy is to be tested empirically and refined through counterexample. Additional roles---strategist, analyst, designer, executor---are interpreted here as differentiated instantiations of the generative function, distinguished by domain of operation rather than by epistemic function.

\textbf{Corollary 1a (proposed).} Under this proposal, the Synthesis function is best read not as a third role but as an emergent property of the Generator/Validator dyad. The agent that produced the thesis position performs synthesis after receiving adversarial feedback, integrating the valid content of the challenge into a refined output. On this reading, synthesis is a phase of the Generator's operation, not a separate epistemic primitive.

\textbf{Corollary 1b (proposed).} On this proposal, the adversarial role is not decomposable into multiple non-adversarial roles. The mode collapse documented in single-generator GAN systems and the conformity effects documented in cooperative multi-agent LLM systems support the reading that adversarial challenge cannot be produced by cooperation among generative agents, regardless of their number \citep{ghosh2018}. Cooperative diversity produces varied generative output; it does not produce an agent with an incentive structure to falsify the outputs of its peers. Adversarial function, on this reading, requires a structural commitment to challenge, not merely a different generative domain.

\subsection{RGAD Tree Structure}
\textbf{Definition (RGAD Tree).} A Recursive Generative-Adversarial Delegation tree is a rooted directed tree $T = (V,E)$ where:
\begin{itemize}
  \item Each node $v \in V$ is an agent (human or AI).
  \item The root node is the terminal human Validator $H$.
  \item Each non-root node $v$ has a parent $\mathrm{parent}(v)$ to whom it serves as Generator.
  \item Each non-leaf node $v$ spawns a Generator child $v_G$ and a Validator child $v_V$ to resolve a sub-problem delegated by $\mathrm{parent}(v)$.
  \item The recursion terminates at leaf nodes, which resolve their assigned sub-problem without further delegation.
\end{itemize}

Three node types are distinguished by structural position:
\begin{itemize}
  \item Root node $H$: the terminal human Validator. A pure Validator with no parent; its discriminative judgment is not itself subject to validation within the system.
  \item Leaf nodes: pure Validators or pure Generators at the base of the recursion. They resolve atomic sub-problems without spawning further children.
  \item Internal non-root nodes: Validator-Orchestrators, occupying both roles simultaneously. Each internal node is a Generator relative to its parent---it produces a synthesized output upward---and a Validator-Orchestrator relative to its children---it discriminates between the outputs of its sub-agents, manages their dialectical exchange, and synthesizes the result into its upward contribution.
\end{itemize}

This dual occupancy is consistent with the manager-agent role in production multi-agent frameworks \citep{hong2024metagpt}: a ``strategist'' or ``manager'' agent is not merely a generator; it is simultaneously the Validator of the outputs produced by the execution agents it supervises. RGAD's contribution at this level is not the discovery of dual-occupancy---which is already widely deployed---but the explicit framing of the same agent's two roles as instances of a single epistemic primitive operating in two directions, which helps explain why human organizational metaphors map intuitively onto multi-agent architectures \citep{lee2026}. The conditions under which a Validator-Orchestrator declares a sub-problem resolved---which classes of stopping criterion are reliable, and how these bound safe autonomous depth---are analyzed in Appendix A.

\subsection{Human as Terminal Validator}
The root node $H$ has a unique structural property: it is a pure Validator with no parent. Its discriminative judgment is not subject to further validation within the system. This non-delegability is structural: the stopping criterion for any sub-tree ultimately derives its validity from the human's goal specification and values, which only the human fully holds.

Replacing the human judge with an AI judge, as in some self-critique or RLAIF settings, does not remove human Validators; it shifts human judgment upstream to the design of the judging model's constitution and objectives. Similarly, in socially consequential settings, the terminal Validator may be a social structure---a court, community, or regulatory body---rather than an individual \citep{rahwan2017}.

\section{Persona Architectures as Cognitive Scaffolding Over RGAD}
With the terminal Validator role established as structurally non-delegable---whether instantiated in an individual human or a social structure---the question becomes how the generative nodes of the RGAD tree should be organized to maximize that Validator's discriminative effectiveness. Persona-based agent architectures offer one class of answer, and we read them as a social-affordance scaffolding layer over the Generator/Validator basis. This section examines the scaffolding mechanism that provides legibility to such architectures and the conditions under which it is necessary.

A common pattern recurs across recent persona-based agent architectures: structural differentiation of generative roles paired with one or a few structurally adversarial Validators. The pattern appears across application domains. In creative problem-solving, persona-split LLMs that distinguish divergent (idea-generating) from convergent (idea-evaluating) modes give the human a clear scaffold for when to expand versus narrow the solution space \citep{epperson2025}. In decision support, multi-persona frameworks deliberately introduce conflicting personas---advocate, skeptic, devil's advocate---so that the human encounters value tensions explicitly rather than implicitly, treating cognitive friction as a design feature \citep{kumar2025}. In human--robot interaction, multi-agent systems with personality-differentiated robots and central coordination improve users' ability to differentiate agents while preserving overall coherence \citep{alami2026}. Persona-based evaluator agents and personalized long-horizon assistants further demonstrate the pattern at evaluation and action layers \citep{li2025, zhang2025}. The Agentic Social Affordance Framework (ASAF) makes the same pattern fully explicit at the system-architecture level with a five-pillar formulation \citep{lee2026}: four pillars (Strategy, Intelligence, Aesthetics, Execution) as differentiated Generators, plus a fifth (Adversarial Audit) structurally committed to challenge. No single persona vocabulary is canonical; what is canonical is that persona architectures of this general shape map onto a Generator-pole differentiation under one or a few structurally adversarial Validators. Read in RGAD terms, this recurring pattern is analogous to a multi-generator GAN: several generators producing diverse outputs under one adversarial pole.

Social legibility---not epistemic function---is what differentiates these Generator-role pillars from a human cognitive perspective. Persona-based designs treat agent identity as a Social Affordance layer: a collaboration interface that structures how human operators perceive and engage each agent. From the RGAD perspective, this implements a scaffolding function that maps abstract Generator/Validator roles onto socially familiar archetypes, allowing human operators to leverage pre-existing social schemas rather than building agent-specific mental models from scratch.

When a human Validator must discriminate among the outputs of many generators, they face a cognitive navigation problem: what criteria apply to each generator, what level of scrutiny is appropriate, what kind of input is expected? For small $N$ (e.g., 2--3 agents), this can be managed procedurally. For larger $N$ (e.g., ten or more specialized agents), the overhead of tracking distinct procedural profiles exceeds working-memory capacity, and discrimination quality degrades.

Social-affordance personas solve this by offloading profile-tracking onto social schemas: a user approaching an ``Inquisitor''-style agent instinctively prepares for adversarial review, supplies more evidence, and applies higher skepticism to outputs without needing to recall a capability list. The persona does cognitive work by activating relevant scripts.

A scaling threshold appears in this picture, at which persona scaffolding shifts from preference to structural necessity. Empirically and theoretically, the threshold aligns with cognitive limits such as Cowan's ``magical number four'' for active items in working memory: beyond roughly four distinct agents, unaided procedural tracking becomes unreliable. At that point persona scaffolding becomes a design requirement rather than a cosmetic choice.

Existing persona vocabularies---divergent/convergent thinker in creativity scaffolding \citep{epperson2025}; advocate/skeptic/devil's-advocate in decision support \citep{kumar2025}; Chief of Staff, Architect, Inquisitor in ASAF \citep{lee2026}---are drawn from organizational and cultural schemas that are cognitively cheap to activate but historically contingent. As human--agent collaboration becomes routine, new schemas specific to this interaction modality are likely to emerge, better aligned with the Generator/Validator structure. Persona scaffolding is therefore best understood as transitional: optimal now, progressively less necessary as native schemas develop. The individual-scale version of this claim is formalized as Hypothesis H4; the civilizational-scale version is a longer-run extrapolation.

\section{The Civilizational Dimension}

\subsection{The Generator-Validator Coupling at the Apex}
The coupling thesis sits in a longer empirical and philosophical tradition than the recent AI-focused literature alone. Goodall's Theory of Expert Leadership documents, across cross-domain empirical studies, that organizations led by people who have done the core work themselves outperform those led by professional managers \citep{goodall2011}. The mechanism Goodall posits---that expert leaders bring ``deep intuitive knowledge''---is what Polanyi's tacit-knowledge tradition formalizes philosophically \citep{polanyi1966}: genuine expertise rests on embodied knowledge that cannot be fully articulated, and therefore cannot be transferred to a validator who has not engaged in the underlying practice. A contemporary analogy may help: the deep intuitive knowledge of experts comprises both education in domain knowledge and first-principles reasoning---analogous to foundation-model pretraining---and years of engaged practice with real-world feedback---analogous to fine-tuning. Without the second, the first does not produce reliable validation. Software engineering specifically has long understood this: empirical research on code review finds that effective discriminative judgment requires sustained engagement with the codebase \citep{bacchelli2013}, and the practitioner debate over whether engineering managers should continue writing code is precisely this coupling problem under a different label.

A growing AI-focused literature reaches the same conclusion through a different door. The Cognitive Integrity Threshold position \citep{lin2026} identifies a comprehension threshold below which oversight becomes ``structurally hollow''---an Empty Oversight regime in which the human can no longer reconstruct or recover system intent. Lin et al. name the underlying phenomenon the \emph{Capability-Comprehension Gap}: a decoupling in which assisted performance improves while the user's internal model of the task deteriorates, so that high apparent competence masks erosion of the comprehension that grounds oversight. Gradual Disempowerment \citep{kulveit2025} documents how incremental AI development can produce permanent, irreversible loss of human influence over the systems civilization depends on; reading this trajectory through RGAD's coupling lens, we extend the analysis to a species-level cognitive selection event in which institutions lose steering capacity---a framing that goes beyond Kulveit et al.'s formulation but is consonant with their concern about gradual erosion of human oversight. The Cognitive Atrophy Paradox literature \citep{kabashkin2025} documents the use-it-or-lose-it dynamics by which over-delegation reduces analytical autonomy. AI deskilling has been argued in this journal to be not an individual but a structural problem \citep{ferdman2025}. And practitioner discourse from defense and security contexts treats human-in-the-loop oversight as ``dangerously misleading'' when operator competence is not actively maintained \citep{dickerson2026, hof2026}. The phenomenon---an apex node that nominally retains authority while losing the competence to exercise it---is well-attested.

What these literatures describe well at the level of phenomenon, they do not fully unify at the level of mechanism. RGAD's contribution at this point is to articulate what we will call the \emph{Generator-Validator coupling thesis}: at any single agent, validation competence is functionally inseparable from sustained generative engagement, because the discriminative judgments a Validator must make depend on the active generative activity that knows what a good candidate looks like. A reviewer who is not actively generating cannot reliably tell a plausible-but-wrong output from a correct one. We propose this thesis as a structural-causal generalization of the expert-leadership and tacit-knowledge traditions, extended to AI delegation in the new cost regime: where generative AI has collapsed the cost of generation by orders of magnitude while leaving validation roughly unchanged, the asymmetric-delegation pattern becomes more attractive than ever, and the coupling-thesis predictions about apex deskilling become correspondingly more acute. Corollary 1a already implies the underlying claim: synthesis is a phase of the Generator's operation, and the Validator that fails to engage with that phase loses access to the criteria by which adversarial challenge is mounted. The two roles are not separable skills coexisting in one agent; they are aspects of the same engaged cognitive activity.

The coupling thesis has a sharp consequence for asymmetric delegation. When the human terminal Validator delegates the Generator role to AI sub-agents while retaining the Validator role nominally, the Validator role does not remain intact. It atrophies in lockstep with the abdicated Generator role---not because humans are inattentive but because the discriminative capacity that constitutes the Validator was always built on, and continuously refreshed by, exercising the Generator. Asymmetric delegation severs the coupling at the precise node where validity terminates, and once severed, the coupling cannot be restored merely by re-asserting Validator authority.

This makes the apex node a single point of failure for the recursive tree, and identifies the specific mechanism: the internal Generator-Validator coupling at one agent is what makes a competent Validator possible at all, and that coupling is not replicable elsewhere in the architecture. Every other node has its work validated by the layer above; only the apex does not. Where the existing literature names the failure (Empty Oversight, gradual disempowerment, structural deskilling), the coupling thesis names the cause.

The thesis predicts a specific empirical signature: AI-assisted humans should lose discriminative skills---debugging, error-detection, plausibility-judgment---faster and more pronouncedly than they lose other forms of competence, because the discriminative skills are the ones constituted by the abdicated Generator engagement. Shen and Tamkin's RCT \citep{shen2026} provides exactly this signature. AI-assisted participants did not just score lower on a knowledge quiz; they showed \emph{particularly pronounced deficits in debugging skill}---the discriminative competence---despite no significant time savings. The qualitative split between users who scaffolded their own reasoning (preserving the Generator role) and users who delegated generation (abdicating it) maps directly onto the coupling thesis: the latter group lost discriminative capacity because they lost the generative engagement that produced it.

The Shen and Tamkin qualitative analysis makes vivid what engagement looks like in practice. The decisive factor in their data was not whether AI was used but \emph{how}: among the AI-interaction patterns they identified, those involving sustained cognitive engagement---asking conceptual questions, reconstructing AI output, treating AI as a scaffold rather than a substitute---preserved learning outcomes, while patterns of full delegation (handing code generation and debugging to AI) surrendered the cognitive activity that produced discriminative skill. The distinction between latent capacity (which may persist for some time after engagement ceases) and competence-as-performance (which requires continuous refreshing through engagement) is what makes asymmetric delegation so corrosive: it can preserve apparent capacity in the short term while removing the engagement that maintains capacity-as-competence over time. The coupling at the apex is, in this dynamic sense, the property of the \emph{engaged human mind}: the cognitive activity of sustained generative participation.

The critical-infrastructure framing follows from this structural argument rather than grounding it---and identifies, in the light of the engagement finding, what the actual infrastructure is. The critical infrastructure of the recursive tree is not the human in some abstract sense (authority-bearer, regulatory agent, biological substrate) but the engaged human mind at the apex: the cognitive activity that constitutes sustained Generator-Validator coupling. The apex node is a single point of failure not because humans are externally important but because the engaged human mind at that one location is what makes validity throughout the system possible. Every node lower in the tree has its work validated by the layer above; only the apex does not, and only the engaged-mind form of the apex sustains the coupling on which apex validation depends. Practitioner concerns at the highest levels of industry---SentinelOne CEO Tomer Weingarten's warning that ``we're deploying systems at global scale that no human understands'' and his framing of human oversight as ``the defining challenge of our lifetime'' \citep{hof2026}---identify the same structural fact in everyday language.

This relocation of the critical infrastructure---from human authority to engaged cognitive activity---has practical consequences for the design of oversight regimes. Interventions that preserve formal authority while permitting engagement to atrophy (signed-off review, rubber-stamp committees, dashboard monitoring of AI outputs) will not preserve the apex coupling and will produce Empty Oversight in the technical sense \citep{lin2026}. Interventions that preserve engagement (active reconstruction of AI reasoning, scheduled rotation through unaided generative work, tutorial-style use of AI in which the human treats the AI as a scaffold for their own thinking) will preserve the coupling even where formal authority is partly delegated downstream. The design question is therefore not how to keep humans \emph{in the loop} but how to keep human minds \emph{engaged in the generative activity} that constitutes the loop.

\subsection{Civilizational Trajectory}
The coupling thesis has implications that extend beyond individual oversight. Civilization can continue to evolve through recursive AI delegation only insofar as the aggregate human Validator infrastructure at the apex can bear the load. The depth and reliability of RGAD trees are bounded by aggregate human discriminative capacity, and that capacity is in turn bounded by preserved generative engagement. A civilization that deepens its AI delegation structures while permitting asymmetric delegation at the apex---generation outsourced, validation nominally retained---is building taller structures on a foundation that is structurally, not contingently, weakening \citep{kulveit2025}.

Historical precedents illustrate a different pattern. Writing externalized memory, freeing cognitive resources; mathematical notation externalized computation; institutions externalized validation into durable social forms. In each case, generative load was offloaded while human comparative advantage concentrated in higher-order generative-discriminative work \citep{wef2026, hogberg2026}. The current transition is novel in that it permits asymmetric delegation of generation specifically: the very competence that grounds discrimination at the apex. Where prior offloading concentrated humans in higher-order generative-discriminative work, present AI delegation can dissolve that work entirely if the Generator-Validator coupling at the apex is not deliberately preserved.

The prediction that the marginal benefit of persona scaffolding diminishes for experienced users (formalized later as Hypothesis H4) is thus not only a measure of user competence; it is also a warning. If persona scaffolding decays too quickly---before native cognitive schemas for multi-agent navigation have formed---human Validators may find themselves in complex delegation trees without either borrowed or native support, precisely when their discriminative judgments are most critical.

\section{Testable Hypotheses}
To move RGAD from a descriptive proposal toward an empirically grounded model, this paper sketches five hypotheses. H1, H4, and H5 are novel to RGAD; H2 and H3 are consonant with empirical predictions already articulated in the persona-scaffolding literature (most explicitly by ASAF \citep{lee2026}) and are restated here in RGAD terms. H5 is the falsifiable form of the Generator-Validator coupling thesis introduced in Section 5.1.

\textbf{H1 (Minimal Basis Sufficiency).} Under equal computational budgets, a two-agent system (one Generator, one Validator) is predicted to achieve comparable accuracy and self-correction to an N-agent system with multiple differentiated Generator roles, when measured on task performance. The advantage of additional persona-based roles is predicted to appear primarily in human interaction quality (e.g., ease of navigation, perceived legibility), not in raw output accuracy.

\textbf{H2 (Identity Signaling and Interaction Framing).} Explicit agent personas and roles will improve human operators' ability to direct tasks to appropriate agents, calibrate oversight, and predict agent behavior, relative to structurally identical systems without persona cues. This is consonant with the Identity Signaling hypothesis in the persona-scaffolding literature.

\textbf{H3 (Behavioral Priming and Collaborative Governance).} Socially differentiated agents (e.g., ``Inquisitor''-style versus ``Architect''-style) will prime humans to provide richer context, anticipate adversarial challenge, and maintain critical distance, reducing over-trust and blind acceptance compared to undifferentiated agents. This restates Behavioral Priming and Collaborative Governance claims from the persona-scaffolding literature.

\textbf{H4 (Scaffolding Decay).} Users with extensive practical exposure to a multi-agent architecture---with the operational threshold for ``extensive'' to be specified by the experimental design---will show diminishing marginal benefits from persona scaffolding on task performance and oversight calibration, compared to novices. This tests the prediction that persona scaffolding is transitional: necessary for novices at scale, but less critical as native interaction schemas develop.

\textbf{H5 (Generator-Validator Coupling).} Under conditions of sustained Generator-role delegation to AI within a domain, human discriminative competence on tasks within that domain (debugging, error-detection, plausibility-judgment, output-quality discrimination) will degrade at a rate exceeding that of other competences during the same delegation period. The mechanism predicted is structural rather than motivational: discriminative competence is constituted by---not independent of---sustained generative engagement, so abdicating the Generator role specifically undermines the Validator role. Falsification would consist of evidence that humans can sustain discriminative performance on a domain over extended Generator-delegation periods without other compensating engagement (e.g., training, practice on adjacent generative tasks). The Shen and Tamkin RCT \citep{shen2026}, with its disproportionate debugging-skill deficit, is consistent with H5 and motivates more direct longitudinal tests.

\section{Limitations}
This work has several limitations.

First, the minimal-basis claim is presented as a proposal motivated by convergence across GANs, dialectics, and debate frameworks; we do not offer a formal proof of minimality. Convergence across three traditions is suggestive, not conclusive. A genuine minimality result would require a fully formalized task model and role-space analysis, which is beyond this paper's scope.

Second, mapping persona architectures such as ASAF's five pillars onto RGAD's Generator half-space plus a structurally adversarial Validator may obscure structural differences between roles (e.g., strategy versus execution). The abstraction treats roles primarily by their epistemic function (generation versus validation), simplifying nuanced differences in local validation.

Third, the civilizational trajectory and critical-infrastructure framing are speculative projections grounded in historical analogy and limited contemporary data. Longitudinal evidence on how professions and institutions reconfigure around RGAD-like structures does not yet exist.

Fourth, this paper does not present new empirical studies run by the author; instead it integrates existing experiments (e.g., Shen and Tamkin's coding-skill RCT) and persona-based design work in the related literature as supporting evidence. The hypotheses in Section 6 remain to be directly tested.

Fifth, RGAD is positioned as both a normative and descriptive proposal, but many deployed systems do not yet conform to its structure. A systematic survey of architectures to empirically assess RGAD as a common pattern or best practice is left to future work.

Finally, of the five hypotheses, H1, H4, and H5 are novel to RGAD; H2 and H3 are restatements of empirical predictions already present in the persona-scaffolding literature. H5 (the Generator-Validator coupling) is currently supported by Shen and Tamkin's debugging-specific deficit pattern but has not been directly tested in a longitudinal design isolating Generator delegation as the independent variable; this paper's empirical originality lies in the reframing and integration of existing evidence rather than in new experiments.

\section{Conclusion}
The contribution of this paper is best characterized as interdisciplinary synthesis rather than structural discovery. Recursive delegation trees with humans at the apex, the generator/validator dyad as a coordination primitive, and the dual-occupancy of internal manager-agents are each well-established in iterated amplification, hierarchical multi-agent system taxonomies, and production manager-worker frameworks \citep{christiano2018ida, christiano2016hch, moore2025, hong2024metagpt}. The structural-causal claim we develop---the Generator-Validator coupling thesis---is itself a generalization of older empirical and philosophical traditions: Goodall's Theory of Expert Leadership \citep{goodall2011} and Polanyi's tacit-knowledge framework \citep{polanyi1966}, with software-engineering-specific empirical grounding in the code-review literature \citep{bacchelli2013}. What we have offered is an epistemically-explicit reading of the recursive-delegation pattern that connects it to the GAN, Hegelian, AI-debate, actor-critic, and human--machine-communication traditions through a unified terminological treatment, and that extends the long-standing insight that competent validation requires sustained engagement with the underlying practice into the new domain of AI delegation under collapsed generation costs.

Within that synthesis lens, RGAD proposes the Generator/Validator dyad as a candidate organizing primitive, models each internal node as simultaneously a Generator to its parent and a Validator-Orchestrator to its children, and---through the Generator-Validator coupling thesis---identifies the engaged human mind at the apex of the tree as the critical infrastructure on which validity throughout the system depends. Critical infrastructure here is not the human as authority-bearer or regulatory agent but the cognitive activity of sustained generative participation, because discriminative competence is constituted by, not independent of, that participation. Reading persona-based architectures---spanning creative collaboration, decision support, human--robot interaction, and explicit system-level frameworks such as ASAF---as social-affordance scaffolding layers over this basis clarifies the role of designed personas and agent identities: they are cognitive scaffolds that make abstract roles legible at scale, structurally necessary when the number of agents exceeds unaided working-memory limits but expected to diminish in necessity as users develop native schemas for human--agent collaboration.

RGAD therefore contributes a unifying conceptual lens rather than a new architecture. It invites further empirical work to test the hypotheses set out in Section 6 and practical experimentation in system design to realize its normative recommendations.

\appendix

\section{Termination Control in Recursive Delegation}
Recursive delegation requires a stopping criterion at each node: when does a Validator-Orchestrator declare a sub-problem resolved and synthesize an output upward? RGAD distinguishes three classes of termination control, ordered by reliability.

\subsection{Axiomatic Reduction}
The most reliable termination occurs when a sub-problem reduces to an external ground truth not generated by the system. Examples include formal proof checking, deterministic unit tests, database queries, and external tools with verifiable behavior. In these cases, termination is grounded in an oracle external to the generative loop.

Leaf nodes that can reach such oracles can complete their work autonomously; internal nodes can terminate when all relevant sub-problems reduce to oracles. Autonomous depth is therefore limited by oracle availability.

\subsection{Convergence Threshold}
When axiomatic reduction is unavailable, internal Validators may apply convergence-based criteria: a Generator's output is accepted if it exceeds a confidence threshold or further iterations no longer produce meaningful improvements. Constitutional AI exemplifies this pattern: self-critique continues until revisions stabilize under a normative constitution \citep{bai2022}.

The vulnerability is confident-but-wrong consensus: multiple cooperative agents can converge on a plausible but incorrect answer without detecting the error. This is the adversarial analog of mode collapse in GANs \citep{goodfellow2014}. External Validators, especially human ones or structurally adversarial agents, remain necessary safeguards.

\subsection{Computational Budget Controls}
The weakest but most universal termination is a hard resource limit: maximum depth, token count, tool calls, or wall-clock time. LLM systems expose such controls as reasoning effort or maximum turns. Budget-based termination is exogenous: recursion stops because a limit is hit, not because the problem is solved. The system returns its best effort so far, analogous to an anytime algorithm.

The practical implication is that safe delegation depth is a function of which termination class is available:

\begin{center}
\begin{tabular}{|l|l|l|l|}
\hline
Termination Class & Signal & Reliability & Max Safe Autonomous Depth \\
\hline
Axiomatic Reduction & External oracle & Highest & Deep -- limited by oracle availability \\
Convergence Threshold & Internal agreement & Medium & Moderate -- conformity risk bounds depth \\
Computational Budget & Resource ceiling & Lowest & Shallow -- human Validator required at exhaustion \\
\hline
\end{tabular}
\end{center}

\section*{Acknowledgments}
This paper was developed in extended dialogue with Claude (Anthropic), accessed via TrendAI's Claude for Enterprise license, used as a research and editorial scaffold rather than as a substitute for the author's own thinking. The collaboration consisted of iterative cycles of conceptual proposal and refinement by the author, with prior-art search, citation verification, structural critique, and prose drafting performed by the AI assistant under the author's editorial control. While TrendAI provides the licensed AI tool, the work was not commissioned by, reviewed by, or written on behalf of TrendAI or Trend Micro Inc.

In retrospect, the working pattern instantiates the engagement-preserving mode the paper identifies in Section 5.1: the author retained the Generator role---originating the framework, the Generator-Validator coupling thesis, the engaged-mind framing, and all substantive judgments---while delegating literature retrieval, characterization of prior work, and drafting to the AI assistant. The author also acknowledges the limits of this pattern: while the most load-bearing citations were independently verified end-to-end, not every cited work was checked exhaustively, which is itself an instance of the asymmetric-delegation risk the paper analyzes. This residual epistemic dependency on the AI assistant's retrieval and characterization is disclosed here in the spirit of the framework rather than concealed.

\section*{Statements and Declarations}

\paragraph{Funding.}
The author received no specific funding for this work.

\paragraph{AI assistance.}
The author used Claude (Anthropic), accessed via TrendAI's Claude for Enterprise license, for editorial assistance during manuscript preparation, including prior-art search, citation verification, structural critique, and prose drafting and refinement. While TrendAI provides the licensed AI tool, the work was not commissioned by, reviewed by, or written on behalf of TrendAI or Trend Micro Inc. The conceptual framework, the Generator-Validator coupling thesis, and all substantive arguments are the author's own. The author reviewed all AI outputs, made independent judgments on what to accept, modify, or reject, and bears full responsibility for the content of this publication. A more reflective account of the working model used appears in the Acknowledgments.

\paragraph{Competing interests.}
The author is a Group Product Manager at TrendAI, Trend Micro Inc. TrendAI develops and markets commercial products in AI-agent security---specifically, AI-agent-powered security tools whose deployment involves agentic AI operating under human oversight. The arguments in this paper, particularly those concerning the need for sustained human generative engagement in oversight of AI agents and the risks of asymmetric delegation, are relevant to the design and positioning of such products. The author has therefore disclosed this potential competing interest. The framework presented here is offered as scholarly analysis and was not commissioned by, reviewed by, or written on behalf of TrendAI or Trend Micro Inc.

\paragraph{Data availability.}
This is a conceptual paper that does not generate or analyze new datasets. All empirical evidence cited is publicly available through the referenced sources.

\paragraph{Author contributions.}
Sole author. The author conceived, drafted, and revised the work.

\paragraph{Ethical approval.}
Not applicable; the paper does not report studies involving human or animal subjects.

\bibliographystyle{plainnat}
\bibliography{references}

\end{document}
